% =============================================================================
% A Bayesian Statistical Framework for Detecting Engine-Assisted Play
% in Online Gomoku Competitions
% Version 3 — Revised with Skill Stratification, Multi-Game Profile Analysis,
%              Sample Size Theory, Experimental Observations,
%              and Opponent-Adjusted Complexity
% =============================================================================
\documentclass[11pt,a4paper]{article}

% --- Packages ---
\usepackage[utf8]{inputenc}
\usepackage[T1]{fontenc}
\usepackage{lmodern}
\usepackage{amsmath,amssymb,amsthm}
\usepackage{graphicx}
\usepackage{booktabs}
\usepackage{hyperref}
\usepackage[margin=2.5cm]{geometry}
\usepackage{algorithm}
\usepackage{algpseudocode}
\usepackage{enumitem}
\usepackage{caption}
\usepackage{subcaption}
\usepackage{xcolor}
\usepackage{float}
\usepackage{natbib}

% --- Theorem environments ---
\newtheorem{definition}{Definition}
\newtheorem{proposition}{Proposition}
\newtheorem{theorem}{Theorem}
\newtheorem{lemma}{Lemma}
\newtheorem{remark}{Remark}

% --- Custom commands ---
\newcommand{\E}{\mathbb{E}}
\newcommand{\Prob}{\mathbb{P}}
\newcommand{\R}{\mathbb{R}}
\DeclareMathOperator*{\argmax}{arg\,max}
\DeclareMathOperator*{\argmin}{arg\,min}

% =============================================================================
\title{
    A Bayesian Statistical Framework for Detecting Engine-Assisted Play\\
    in Online Gomoku Competitions
}

\author{
    Nguyen Cong Minh \\
    \texttt{nguyenminhfptu@gmail.com}
}

\date{\today}

% =============================================================================
\begin{document}

\maketitle

% =============================================================================
% ABSTRACT
% =============================================================================
\begin{abstract}
The integrity of online board game competitions is increasingly threatened by
players who covertly use artificial intelligence engines during live matches.
This paper presents a rigorous Bayesian statistical framework for detecting
engine-assisted play in Gomoku (Five-in-a-Row), a combinatorial strategy game.
Our approach models the \emph{winrate delta}---the difference between the
engine-optimal move evaluation and the player's actual move evaluation---as a
random variable drawn from parametric distributions that differ between
legitimate human players and engine-assisted players. We develop the framework
within a general family of non-negative distributions (Exponential, Gamma,
Weibull), and introduce a \emph{two-component mixture model} to capture the
bimodal nature of human play, with a Dirichlet Process extension for
automatic component selection. A \emph{context-aware complexity metric} weights
evidence by positional difficulty, and a \emph{tempered likelihood} mechanism
prevents trivial positions from distorting posterior estimates. The framework
performs sequential Bayesian updates over the course of a game, integrating
a Boltzmann choice model, information-theoretic divergence measures, and
temporal switch-point detection to identify mid-game engine activation.
We further introduce \emph{player skill stratification} to address the
false-positive problem arising from strong legitimate players, a formal
\emph{multi-game profile analysis} grounded in confidence-interval theory
and hypothesis testing for population-level inference, and an
\emph{opponent-adjusted complexity} mechanism that corrects for positions
made artificially easy by a weak opponent---a previously unaddressed source
of false-positive inflation. A formal likelihood ratio test grounded in
Neyman--Pearson theory provides false-positive rate control. We describe
the complete mathematical formulation, the system architecture for engine
integration, and the classification decision boundaries, and report
preliminary experimental observations that motivate several design choices.

\medskip
\noindent
\textbf{Keywords:} Anti-cheating, Bayesian inference, Gomoku, engine-assisted
play detection, online gaming integrity, mixture model, Gamma distribution,
Weibull distribution, Dirichlet process, information theory, temporal analysis,
skill stratification, profile analysis, opponent-adjusted complexity.
\end{abstract}

% =============================================================================
% 1. INTRODUCTION
% =============================================================================
\section{Introduction}
\label{sec:introduction}

Online board game platforms have experienced significant growth, enabling
players worldwide to compete in real-time tournaments. However, the remote and
unsupervised nature of online play creates an avenue for a specific form of
cheating: covertly consulting an AI engine during a live match. In Gomoku
(also known as Five-in-a-Row or Renju in its professional variant), modern
engines achieve superhuman playing strength, making engine-assisted play both
highly effective and difficult to distinguish from expert human play through
casual observation alone.

Existing anti-cheating measures in online gaming have largely relied on
behavioral heuristics---such as timing patterns, mouse movement tracking,
or tab-switching detection---which are easily circumvented by sophisticated
cheaters. Statistical approaches have been explored in chess
\citep{regan2011intrinsic}, but relatively little formal work has been
published for Gomoku and its variants. Furthermore, many deployed systems
rely on ad-hoc thresholds (e.g., ``if accuracy exceeds 90\%, flag the
player'') without a coherent probabilistic foundation, leading to both
false positives against strong legitimate players and false negatives
against selective cheaters.

A fundamental challenge, which we term the \emph{skill-confounding problem},
arises because a single-population null hypothesis $H_0$ cannot simultaneously
represent beginners, intermediate players, and professionals. A professional
player's move quality may exceed the system's cheater threshold while being
entirely consistent with legitimate play. Preliminary experiments (Section~%
\ref{sec:experiments}) confirm this: a professional player analyzed over
22~moves yielded $\hat{\lambda}_{\text{MLE}} = 113$, far exceeding the default
cheater threshold of $\lambda_{\text{cheater}} = 50$, yet the player was
confirmed to be playing without assistance. This motivates the stratified null
hypothesis developed in Section~\ref{sec:stratification}.

This paper addresses these limitations by developing a complete Bayesian
statistical framework that:

\begin{enumerate}[label=(\roman*)]
    \item Models human and cheater move-quality distributions using parametric
          families grounded in statistical theory.
    \item Incorporates \emph{positional context} through a complexity metric
          that modulates the evidential weight of each move.
    \item Performs online (sequential) Bayesian updates with tempered
          likelihoods, ensuring that trivial positions do not corrupt the
          posterior estimate.
    \item Employs information-theoretic divergence measures and temporal
          change-point detection to identify partial or mid-game engine use.
    \item Stratifies the null hypothesis by player skill level and aggregates
          evidence across multiple games via a formal profile analysis.
    \item Corrects for \emph{opponent-induced easy positions} via an
          opponent-adjusted complexity mechanism that prevents artificially
          inflated posteriors when the analyzed player faces a weak opponent.
\end{enumerate}

\subsection{Contributions}

The main contributions of this work are:

\begin{enumerate}
    \item \textbf{Winrate Delta Model.} We formalize the concept of
          \emph{winrate delta} and establish that the Exponential distribution
          family provides a natural parametric model for move quality under both
          hypotheses, with Gamma and Weibull as empirically motivated
          generalizations.
    \item \textbf{Two-Component Mixture Model.} We propose a mixture of
          Exponentials for the human hypothesis that accounts for the bimodal
          nature of human play (careful moves interspersed with occasional
          blunders), with parameters estimated via the EM algorithm.
    \item \textbf{Player Skill Stratification.} We introduce a stratified null
          hypothesis that conditions on the player's estimated skill tier,
          resolving the false-positive problem for strong legitimate players.
    \item \textbf{Context-Aware Complexity.} We derive a composite complexity
          metric from candidate move evaluations that weights Bayesian evidence
          by positional difficulty, preventing trivial positions from generating
          misleading signals.
    \item \textbf{Tempered Bayesian Updates.} We introduce a tempering mechanism
          where the likelihood exponent equals the positional complexity,
          providing a mathematically principled method for evidence weighting.
    \item \textbf{Multi-Game Profile Analysis.} We develop a formal
          population-level inference framework that aggregates evidence across
          multiple games, with sample size requirements derived from
          confidence-interval theory.
    \item \textbf{Temporal Switch-Point Detection.} We develop a temporal model
          that detects abrupt changes in play quality, identifying possible
          mid-game engine activation.
    \item \textbf{Opponent-Adjusted Complexity.} We identify and formalize the
          \emph{opponent-difficulty confound}---the systematic inflation of
          false-positive rates when the analyzed player faces a weak opponent---
          and introduce the Opponent Difficulty Factor to correct for it.
    \item \textbf{Complete System Design.} We describe the full architecture for
          integrating an external game engine as an evaluation oracle within the
          detection pipeline.
\end{enumerate}

\subsection{Paper Organization}

Section~\ref{sec:related_work} reviews related work. Section~\ref{sec:framework}
presents the mathematical framework in detail. Section~\ref{sec:complexity}
develops the context-aware complexity model. Section~\ref{sec:bayesian}
describes the Bayesian inference engine. Section~\ref{sec:information_theory}
introduces the information-theoretic metrics. Section~\ref{sec:temporal}
covers temporal analysis and switch-point detection. Section~\ref{sec:profile}
develops the multi-game profile analysis framework. Section~%
\ref{sec:architecture} outlines the system architecture.
Section~\ref{sec:classification} defines the classification decision
boundaries. Section~\ref{sec:experiments} reports preliminary experimental
observations. Section~\ref{sec:discussion} discusses strengths, limitations,
and ethical considerations. Section~\ref{sec:conclusion} concludes the paper.


% =============================================================================
% 2. RELATED WORK
% =============================================================================
\section{Related Work}
\label{sec:related_work}

\subsection{Engine-Assisted Play Detection in Chess}

The most extensively studied domain for engine-assisted play detection is chess.
Regan and Haworth~\citep{regan2011intrinsic} proposed an intrinsic performance
rating system that estimates a player's Elo rating from their moves alone,
comparing observed play quality to engine evaluations. Their approach uses a
logistic model to predict the probability that a player of a given rating would
select each candidate move. Significant deviations between the predicted rating
and the player's registered rating are flagged as suspicious.

The \emph{Let's Check} system used by ChessBase and similar platforms
\citep{guid2006computer} compares player moves against deep engine analysis,
computing aggregate statistics such as average centipawn loss and correlation
with engine top choices. While effective as a screening tool, these systems
typically rely on fixed thresholds rather than probabilistic inference.

\subsection{Statistical Methods for Fair Play}

Broader statistical approaches to fair play detection include hypothesis
testing frameworks~\citep{berger1987testing}, where the null hypothesis
is honest play and the alternative is engine assistance. The challenge
lies in controlling the false positive rate---strong human players
naturally exhibit high move quality, and a framework that flags top-level
play as suspicious undermines competitive integrity.

\subsection{Gomoku-Specific Considerations}

Gomoku differs from chess in several important respects that affect anti-cheating
analysis. The game has a narrower decision tree in many positions due to forced
tactical sequences (threats, double-threats), meaning that many ``perfect''
moves are also the obvious human choice. A robust detection framework must
therefore distinguish between positions where engine-optimal play is trivially
findable and positions where it requires deep calculation. This motivates our
development of the context-aware complexity model (Section~\ref{sec:complexity}).


% =============================================================================
% 3. MATHEMATICAL FRAMEWORK
% =============================================================================
\section{Mathematical Framework: The Winrate Delta Model}
\label{sec:framework}

\subsection{Definitions}

We assume access to an external game engine that, given a board position~$s$,
produces a set of $k$ candidate moves $\{m_1, m_2, \ldots, m_k\}$ along with
their associated \emph{winrate evaluations} $\{w_1, w_2, \ldots, w_k\}$, where
$w_i \in [0, 1]$ represents the probability of winning from the resulting
position if move~$m_i$ is played.

\begin{definition}[Winrate Delta]
\label{def:delta}
For a given position, let $w^* = \max_i w_i$ denote the winrate of the
engine's best candidate move, and let $w_p$ denote the winrate of the move
actually played by the player. The \emph{winrate delta} is defined as:
\begin{equation}
    \Delta = w^* - w_p, \quad \Delta \geq 0.
    \label{eq:delta}
\end{equation}
\end{definition}

The winrate delta $\Delta$ captures the \emph{quality loss} incurred by the
player's move choice. A $\Delta$ of zero indicates the player found the
engine-optimal move; a large $\Delta$ indicates a suboptimal choice.

\begin{remark}
All winrates are expressed from the perspective of the player being analyzed.
If the engine reports winrates from the opponent's perspective, the inversion
$w_{\text{player}} = 1 - w_{\text{opponent}}$ must be applied before computing
$\Delta$.
\end{remark}

\subsection{Distributional Family for the Winrate Delta}
\label{sec:dist_family}

Because $\Delta \geq 0$ by definition, the appropriate parametric models
belong to the family of non-negative continuous distributions. We consider
three candidates in order of increasing flexibility.

\subsubsection{Exponential Distribution (Baseline)}

The simplest model assumes $\Delta \sim \text{Exponential}(\lambda)$:
\begin{equation}
    f(\Delta \mid \lambda) = \lambda \, e^{-\lambda \Delta}, \quad \Delta \geq 0.
    \label{eq:exponential_pdf}
\end{equation}
Its memoryless property implies that the magnitude of a suboptimal choice does
not depend on previous errors conditional on the current position. However,
this assumption may be violated when players exhibit psychological momentum
(``tilt'' or ``flow state''), producing serial correlation in error magnitudes.

\subsubsection{Gamma Distribution (Shape-Rate Generalization)}

\begin{equation}
    f(\Delta \mid \alpha, \beta) = \frac{\beta^{\alpha}\,
    \Delta^{\alpha - 1}\, e^{-\beta\Delta}}{\Gamma(\alpha)},
    \quad \Delta \geq 0.
    \label{eq:gamma_pdf}
\end{equation}
When $\alpha = 1$, the Gamma reduces to $\text{Exponential}(\beta)$.
For $\alpha > 1$, the density is unimodal with mode at
$(\alpha - 1)/\beta$, enabling modeling of delta distributions that
concentrate around a typical error magnitude. The mean and variance are
$\E[\Delta] = \alpha/\beta$ and $\text{Var}(\Delta) = \alpha/\beta^2$.

\subsubsection{Weibull Distribution (Tail Flexibility)}

\begin{equation}
    f(\Delta \mid k, \lambda) = \frac{k}{\lambda}
    \left(\frac{\Delta}{\lambda}\right)^{k-1}
    \exp\!\left[-\left(\frac{\Delta}{\lambda}\right)^{k}\right],
    \quad \Delta \geq 0.
    \label{eq:weibull_pdf}
\end{equation}
The shape parameter $k$ controls tail weight: $k < 1$ yields a heavier tail
than the Exponential (suitable for humans prone to rare severe blunders), while
$k > 1$ yields a lighter tail (suitable for cheater modeling). Preliminary
experiments consistently select Weibull as the best-fit model by AIC/BIC
(Section~\ref{sec:experiments}), though all models are rejected by
goodness-of-fit tests at small sample sizes ($n \leq 30$), underscoring the
need for multi-game aggregation.

\subsubsection{Model Selection via Information Criteria}
\label{sec:model_selection}

\begin{align}
    \text{AIC} &= 2k - 2\ln \hat{L}, \label{eq:aic}\\
    \text{BIC} &= k \ln n - 2\ln \hat{L}. \label{eq:bic}
\end{align}
Lower values indicate better model--data fit after penalizing complexity.
Goodness-of-fit tests (Kolmogorov--Smirnov or Anderson--Darling) should
additionally be applied to validate distributional assumptions before deployment.

\subsection{Hypothesis Structure and Player Skill Stratification}
\label{sec:stratification}

A critical deficiency of single-population null hypothesis models is the
\emph{skill-confounding problem}: a professional player's move quality
may be statistically indistinguishable from engine-assisted play under a
model calibrated for average-strength players. We address this by
introducing a \emph{stratified null hypothesis}.

\subsubsection{Skill Tiers}

We partition the player population into skill tiers $\mathcal{T} =
\{T_1, T_2, T_3, T_4\}$, each characterized by a distinct null
distribution parameter:

\begin{table}[H]
\centering
\caption{Proposed skill tier parameterization for the Exponential baseline.
Values should be calibrated empirically from labeled reference data.}
\label{tab:tiers}
\begin{tabular}{llll}
\toprule
Tier & Description & $\lambda_{H_0}$ (indicative) & $\E[\Delta]$ \\
\midrule
$T_1$ & Beginner       & $3$--$8$   & $12$--$33\%$ \\
$T_2$ & Intermediate   & $10$--$25$ & $4$--$10\%$  \\
$T_3$ & Advanced       & $30$--$60$ & $1.7$--$3.3\%$ \\
$T_4$ & Professional   & $80$--$150$ & $0.7$--$1.3\%$ \\
\hline
$H_1$ & Engine-assisted & $200$--$500$ & $0.2$--$0.5\%$ \\
\bottomrule
\end{tabular}
\end{table}

\begin{remark}
The values in Table~\ref{tab:tiers} are indicative only. Empirical calibration
from large labeled datasets (Section~\ref{sec:em}) is required before
deployment. A key observation from preliminary experiments is that a
professional player's $\hat{\lambda}_{\text{MLE}}$ consistently falls in
the range $90$--$120$ (see Section~\ref{sec:experiments}), firmly within tier
$T_4$, while the current default $\lambda_{\text{cheater}} = 50$ lies
\emph{within} the professional human range---demonstrating the necessity of
stratification.
\end{remark}

\subsubsection{Tier Assignment}

In practice, the player's tier is assigned based on one of three methods,
in decreasing order of reliability:
\begin{enumerate}
    \item \textbf{Profile-estimated tier:} The player's historical
          $\hat{\lambda}$ from confirmed honest games (Section~\ref{sec:profile})
          directly determines tier membership.
    \item \textbf{Registered rating:} The player's tournament rating (Elo or
          equivalent) maps to a tier via calibrated thresholds.
    \item \textbf{Prior tier:} In the absence of historical data, a weakly
          informative prior assigns the player to the most common tier for
          the relevant competitive context.
\end{enumerate}

\subsubsection{Stratified Likelihood}

Under the stratified null hypothesis, the human likelihood in
Equation~\eqref{eq:bayes} is replaced by the tier-conditional likelihood:
\begin{equation}
    P(\Delta \mid H_0) = P(\Delta \mid H_0, T^*),
    \label{eq:stratified_h0}
\end{equation}
where $T^*$ is the player's assigned tier. This ensures that the
Bayesian posterior compares the observed play to the player's
\emph{own peer group} rather than the overall player population.

\subsection{Maximum Likelihood Estimation}

Given a sequence of observed deltas $\Delta_1, \Delta_2, \ldots, \Delta_n$,
the MLE for the Exponential rate parameter is:
\begin{equation}
    \hat{\lambda}_{\text{MLE}} = \frac{n}{\sum_{i=1}^{n} \Delta_i}
    = \frac{1}{\bar{\Delta}}.
    \label{eq:lambda_mle}
\end{equation}

To handle $\Delta_i = 0$, we apply a floor $\epsilon = 0.001$:
\begin{equation}
    \Delta_i^{\text{adj}} = \max(\Delta_i, \, \epsilon).
    \label{eq:delta_adj}
\end{equation}

We define the \emph{Accuracy Index} as:
\begin{equation}
    \text{AccuracyIndex} = \frac{1}{\hat{\lambda}_{\text{MLE}}}
    = \bar{\Delta}^{\text{adj}},
    \label{eq:accuracy_index}
\end{equation}
which directly measures the average winrate loss. Lower values indicate play
closer to engine-optimal.

\subsection{Two-Component Exponential Mixture for Human Play}
\label{sec:mixture}

A single Exponential distribution inadequately captures human play, which
exhibits a bimodal pattern: the majority of moves are sound (small $\Delta$),
but occasional moves are blunders (large $\Delta$). We model the human
delta distribution as:
\begin{equation}
    P(\Delta \mid H_0) = \pi \cdot \lambda_g \, e^{-\lambda_g \Delta}
    + (1 - \pi) \cdot \lambda_b \, e^{-\lambda_b \Delta},
    \label{eq:mixture_pdf}
\end{equation}
where $\pi \in (0, 1)$ is the mixing weight for the ``good play'' component,
$\lambda_g$ is the rate for good plays, and $\lambda_b$ is the rate for
blunders.

The log-sum-exp trick ensures numerical stability:
\begin{align}
    a &= \log \pi + \log \lambda_g - \lambda_g \Delta, \\
    b &= \log(1 - \pi) + \log \lambda_b - \lambda_b \Delta, \\
    \log P(\Delta \mid H_0) &= m + \log\!\left(e^{a-m} + e^{b-m}\right),
\end{align}
where $m = \max(a, b)$.

\subsection{Parameter Estimation via Expectation-Maximization}
\label{sec:em}

The mixture parameters $\Theta = (\pi, \lambda_g, \lambda_b)$ are estimated
from labeled reference data via the EM algorithm~\citep{dempster1977maximum}.

\textbf{E-step:}
\begin{equation}
    \gamma_{i,1} = \frac{\pi \cdot \lambda_g e^{-\lambda_g \Delta_i}}
    {\pi \cdot \lambda_g e^{-\lambda_g \Delta_i}
    + (1-\pi) \cdot \lambda_b e^{-\lambda_b \Delta_i}}, \quad
    \gamma_{i,2} = 1 - \gamma_{i,1}.
    \label{eq:e_step}
\end{equation}

\textbf{M-step:}
\begin{align}
    \pi^{\text{new}} &= \frac{1}{n}\sum_{i=1}^{n} \gamma_{i,1},
    \label{eq:m_step_pi}\\
    \lambda_g^{\text{new}} &= \frac{\sum_i \gamma_{i,1}}
    {\sum_i \gamma_{i,1} \Delta_i}, \quad
    \lambda_b^{\text{new}} = \frac{\sum_i \gamma_{i,2}}
    {\sum_i \gamma_{i,2} \Delta_i}.
    \label{eq:m_step_lambda}
\end{align}

\begin{remark}
When EM is applied to a single-game dataset ($n \leq 30$), the fitted
parameters exhibit high variance and should be treated as exploratory
descriptors rather than calibrated estimates. Reliable EM convergence
requires $n \geq 200$ observations (Section~\ref{sec:sample_size}).
\end{remark}

\subsection{Dirichlet Process Mixture Model}
\label{sec:dpm}

The fixed two-component mixture assumes exactly two behavioral modes.
The \emph{Dirichlet Process Mixture Model} (DPMM)~\citep{ferguson1973bayesian,
neal2000markov} relaxes this assumption:
\begin{equation}
    G \sim \text{DP}(\alpha_0, G_0), \qquad
    \theta_i \mid G \sim G, \qquad
    \Delta_i \mid \theta_i \sim F(\cdot \mid \theta_i).
\end{equation}
The effective number of components is governed by the concentration parameter
$\alpha_0$ and inferred from data, avoiding model selection entirely.

\subsection{Boltzmann Choice Model}
\label{sec:softmax}

We model the player's move selection using a Boltzmann distribution:
\begin{equation}
    P(m_i \mid \{w_j\}, \tau) = \frac{\exp(w_i / \tau)}{\sum_{j=1}^{k}
    \exp(w_j / \tau)},
    \label{eq:softmax}
\end{equation}
where $\tau \to 0$ concentrates mass on the highest-winrate move (engine-like)
and $\tau \to \infty$ produces a uniform distribution (random play).

The temperature is estimated by maximum likelihood:
\begin{equation}
    \hat{\tau} = \argmax_{\tau} \sum_{i=1}^{n} \log P(m_{c_i} \mid
    \{w_j^{(i)}\}, \tau).
\end{equation}

The temperature score $S_{\tau} = 1/\hat{\tau}$ enters the classification
pipeline as an auxiliary signal (Section~\ref{sec:classification}).


% =============================================================================
% 4. CONTEXT-AWARE COMPLEXITY
% =============================================================================
\section{Context-Aware Position Complexity}
\label{sec:complexity}

\subsection{Delta Vector and Key Metrics}

\begin{equation}
    d_i = w_i - w_{\text{prev}}, \quad
    d^* = \max_i d_i, \quad
    \bar{d} = \frac{1}{k}\sum_{i} d_i, \quad
    S = d^* - \min_i d_i.
\end{equation}

\subsection{Complexity Factors}

\begin{definition}[Impact Factor]
$F_{\text{impact}} = \min\!\left(|d^*| / 0.25,\; 1.0\right).$
\end{definition}

\begin{definition}[Variance Factor]
$F_{\text{variance}} = \min\!\left(S / 0.40,\; 1.0\right).$
\end{definition}

\begin{definition}[Criticality Factor]
$F_{\text{criticality}} = (d^* - \bar{d}) / S$ if $S > 0.01$, else $0$,
clamped to $[0,1]$.
\end{definition}

\begin{definition}[Phase Factor]
\begin{equation}
    F_{\text{phase}} = \begin{cases}
        0.3 & \text{if } w^* > 0.85, \\
        0.2 & \text{if } w^* < 0.15, \\
        0.5 + 0.5(1 - 2|w^* - 0.5|) & \text{otherwise}.
    \end{cases}
\end{equation}
\end{definition}

\subsection{Combined Complexity Score}

\begin{equation}
    C = 0.15 F_{\text{impact}} + 0.40 F_{\text{variance}}
    + 0.20 F_{\text{criticality}} + 0.25 F_{\text{phase}}.
    \label{eq:complexity}
\end{equation}

\begin{proposition}[Trivial Position Override]
\label{prop:trivial}
If $w^* > 0.95$ or $w^* < 0.05$ and $F_{\text{variance}} < 0.05$,
then $C = 0.10$.
\end{proposition}

\subsection{Opponent-Adjusted Complexity}
\label{sec:opp_complexity}

The complexity score $C$ defined above depends solely on the \emph{analyzed
player's} candidate move evaluations, and is therefore blind to the state
of the opponent's position. This creates a systematic bias that we term
the \emph{opponent-difficulty confound}.

\begin{definition}[Opponent-Difficulty Confound]
\label{def:opp_confound}
When the analyzed player's opponent is in a severely losing position,
many of the player's candidate moves achieve near-optimal winrate
(because the opponent has no effective defense). A one-sided complexity
model that ignores opponent position assigns non-negligible complexity
$C$ to these positions, incorrectly treating trivially-optimal play as
informative evidence of engine assistance.
\end{definition}

\begin{remark}[Empirical Motivation]
\label{rem:opp_motivation}
In a preliminary game analysis (Section~\ref{sec:exp_opp}), multiple
late-game positions were assigned complexity $C \approx 0.49$ despite the
opponent's best available move having a winrate below $0.3\%$.
The Trivial Position Override (Proposition~\ref{prop:trivial}) failed to
activate because $w^*$ had not yet exceeded the $0.95$ threshold from the
player's perspective, while the opponent was already effectively defeated.
These positions contributed spurious evidence toward $H_1$.
\end{remark}

\subsubsection{Opponent Difficulty Factor}

Let $w^*_{\text{opp}}$ denote the best available winrate for the opponent
at position $t$, expressed from the \emph{opponent's perspective} (i.e.,
already normalized via $w_{\text{opp}} = 1 - w_{\text{player}}$ per
Remark~1). The \emph{Opponent Difficulty Factor} is defined as:

\begin{equation}
    F_{\text{opp}} = \min\!\left(\frac{w^*_{\text{opp}}}{w_{\text{ref}}},\;
    1.0\right),
    \label{eq:f_opp}
\end{equation}
where $w_{\text{ref}} = 0.40$ is a calibration threshold representing the
winrate below which the opponent is no longer considered a meaningful
competitive threat. The factor $F_{\text{opp}} \in [0, 1]$ takes value~$1$
when the opponent is fully competitive and decreases toward~$0$ as the
opponent's position becomes hopeless.

\begin{remark}[Temporal Indexing]
To avoid circular dependency, $w^*_{\text{opp}}$ at position $t$ is
evaluated \emph{after the opponent's most recent move} and \emph{before}
the analyzed player's current move is played. This ensures that we measure
the opponent's strength at the moment the decision is made, not the
consequence of the player's own move.
\end{remark}

\subsubsection{Opponent-Adjusted Complexity Score}

The raw complexity $C$ (Equation~\eqref{eq:complexity}) is adjusted as:

\begin{equation}
    C_{\text{adj}} = C \cdot \max\!\left(F_{\text{opp}},\; \alpha_{\min}\right),
    \label{eq:c_adj}
\end{equation}
where $\alpha_{\min} = 0.10$ ensures that a minimum residual weight is
retained even in extreme cases, preserving the continuity of the evidence
accumulation process.

The final complexity used in all downstream computations (tempered
likelihood, context-aware accuracy) is:
\begin{equation}
    C_{\text{final}} = \begin{cases}
        0.10 & \text{if Trivial Position Override is triggered,} \\
        C_{\text{adj}} & \text{otherwise.}
    \end{cases}
    \label{eq:c_final}
\end{equation}

\begin{remark}[Effect on Tempered Likelihood]
Note that $C_{\text{final}}$ acts as the tempering exponent $\alpha$ in
Equation~\eqref{eq:tempering}. Therefore, when the opponent is nearly
defeated ($F_{\text{opp}} \approx 0$), the tempered likelihood approaches
$P(\Delta \mid H)^{\alpha_{\min}} = P(\Delta \mid H)^{0.10}$, contributing
only minimal evidence to the Bayesian posterior regardless of how small $\Delta$
is. This is the mathematically correct behavior: a perfect move in a
trivially-won position reveals nothing about the player's computational
capability.
\end{remark}

Table~\ref{tab:opp_example} illustrates the effect of the adjustment on
representative positions from the experimental game log.

\begin{table}[H]
\centering
\caption{Effect of Opponent-Adjusted Complexity on selected moves
(experimental game, Black side). $C$ denotes raw complexity;
$F_{\text{opp}}$ denotes the Opponent Difficulty Factor;
$C_{\text{adj}}$ denotes the adjusted complexity.}
\label{tab:opp_example}
\begin{tabular}{ccccccl}
\toprule
Move & $\Delta$ & $C$ & $w^*_{\text{opp}}$ & $F_{\text{opp}}$ &
$C_{\text{adj}}$ & Opp Quality \\
\midrule
\#7  & 0.0\% & 0.56 & 36.9\% & 0.92 & 0.515 & Passive (competitive) \\
\#9  & 3.5\% & 0.46 & 36.4\% & 0.91 & 0.419 & Tactical (competitive) \\
\#11 & 0.0\% & 0.62 & 22.9\% & 0.57 & 0.355 & Tactical (weakening) \\
\#13 & 0.0\% & 0.47 & 0.2\%  & \textbf{0.005} & \textbf{0.047} & Winning (trivial) \\
\#33 & 0.0\% & 0.49 & 0.2\%  & \textbf{0.005} & \textbf{0.049} & Winning (trivial) \\
\#35 & 0.0\% & 0.10 & 0.2\%  & 0.005 & 0.010 & Override applied \\
\bottomrule
\end{tabular}
\end{table}

Without the adjustment, moves \#13 and \#33 contribute $C = 0.47$--$0.49$
worth of evidence toward $H_1$ despite occurring in trivially-decided
positions. With $C_{\text{adj}}$, their evidential contribution drops to
$\approx 0.05$, commensurate with their actual informativeness.

\subsection{Context-Aware Accuracy}

All accuracy and likelihood computations use $C_{\text{final}}$
(Equation~\eqref{eq:c_final}) in place of the raw complexity $C$:
\begin{equation}
    A(\Delta, C_{\text{final}}) = \begin{cases}
        1.0 & \text{if } \Delta \leq 0.02, \\
        \exp\!\left(-5 \Delta (1 - 0.3 C_{\text{final}})\right) &
        \text{otherwise.}
    \end{cases}
\end{equation}

\subsubsection{Strategic Bonus}

\begin{equation}
    A_{\text{final}} = A \cdot (1 + 0.15 C_{\text{opp}}),
    \quad \text{if } A > 0.4 \text{ and } C_{\text{opp}} > 0.1,
\end{equation}
clamped to $[0, 1]$.


% =============================================================================
% 5. BAYESIAN INFERENCE
% =============================================================================
\section{Bayesian Inference Engine}
\label{sec:bayesian}

\subsection{Posterior Computation}

\begin{equation}
    P(H_1 \mid D) = \frac{P(D \mid H_1) \cdot P(H_1)}{P(D \mid H_1) \cdot
    P(H_1) + P(D \mid H_0) \cdot P(H_0)},
    \label{eq:bayes}
\end{equation}
with $P(H_1) = 0.01$ (conservative prior) and $P(H_0) = 1 - P(H_1)$.

In log space:
\begin{align}
    L_0 &= \sum_{i=1}^{n} \log P(\Delta_i \mid H_0), &
    L_1 &= \sum_{i=1}^{n} \log P(\Delta_i \mid H_1),\\
    P(H_1 \mid D) &= \frac{e^{\ell_1 - m}}{e^{\ell_0 - m} + e^{\ell_1 - m}},
\end{align}
where $\ell_j = L_j + \log P(H_j)$ and $m = \max(\ell_0, \ell_1)$.

\subsection{Online (Sequential) Bayesian Updates}

\begin{equation}
    P(H_1 \mid D_{1:n+1}) \propto P(\Delta_{n+1} \mid H_1)^{C_{n+1}} \cdot
    P(H_1 \mid D_{1:n}).
\end{equation}

\subsection{Tempered Likelihood}
\label{sec:tempered}

We raise the likelihood to the power of the positional complexity $C$:
\begin{equation}
    \label{eq:tempering}
    P(\Delta \mid H)^{C} \quad \Longleftrightarrow \quad
    C \cdot \log P(\Delta \mid H).
\end{equation}
When $C = 0$ (trivial position), the tempered likelihood is~$1$, contributing
zero evidence. When $C = 1$, full evidence is contributed. This mechanism
is well-studied in the statistical literature~\citep{grunwald2017inconsistency}.

\begin{algorithm}[H]
\caption{Online Bayesian Update with Tempered Likelihood}
\label{alg:online_update}
\begin{algorithmic}[1]
\Require Current posterior $p$, new delta $\Delta$, complexity $C$
\Ensure Updated posterior $p'$
\If{$C \leq 0$}
    \State \Return $p$
\EndIf
\State $\ell_0 \gets C \cdot \log P(\Delta \mid H_0) + \log(1 - p)$
\State $\ell_1 \gets C \cdot \log P(\Delta \mid H_1) + \log(p)$
\State $m \gets \max(\ell_0, \ell_1)$
\State \Return $e^{\ell_1 - m} / (e^{\ell_0 - m} + e^{\ell_1 - m})$
\end{algorithmic}
\end{algorithm}

\subsection{Confidence Estimation}

\begin{equation}
    \text{confidence} = \min\!\left(1, \frac{n}{n_0}\right)
    \times 2 \left|P(H_1 \mid D) - 0.5\right|,
    \label{eq:confidence}
\end{equation}
where $n_0$ is a data-sufficiency threshold. The value $n_0 = 30$ used in
the original formulation is a lower bound; Section~\ref{sec:sample_size}
derives principled values based on the desired estimation precision. We
recommend $n_0 = 44$ for proportion-based classification and $n_0 = 97$
for $\lambda$-based inference.

\subsection{Formal Likelihood Ratio Test}
\label{sec:lrt}

\begin{equation}
    \Lambda = \frac{P(D \mid H_1)}{P(D \mid H_0)}
    = \exp(L_1 - L_0).
    \label{eq:lrt}
\end{equation}
By the Neyman--Pearson lemma~\citep{neyman1933most}, the test that rejects
$H_0$ when $\Lambda > \eta$ is the most powerful test at significance
level $\alpha$. The critical value $\eta$ is calibrated from a reference
corpus of confirmed human games.

\subsection{Prior Sensitivity Analysis}
\label{sec:sensitivity}

We evaluate $P(H_1 \mid D)$ across
$P(H_1) \in \{0.001, 0.01, 0.05, 0.10\}$. If the posterior exceeds the
cheater threshold for all reasonable priors, the evidence is considered robust.
Sensitivity analysis is standard practice in Bayesian applications with
subjective priors~\citep{berger1987testing}.


% =============================================================================
% 6. INFORMATION THEORY
% =============================================================================
\section{Information-Theoretic Metrics}
\label{sec:information_theory}

\subsection{Shannon Entropy of Move Quality}

\begin{equation}
    H(P) = -\sum_{i} p_i \log_2 p_i,
\end{equation}
where $p_i$ is the empirical probability of the player's delta falling in bin
$i$. Human players exhibit higher entropy (more varied quality); engine-assisted
players exhibit lower entropy (consistently near-optimal).

\subsection{KL Divergence from Engine Policy}

\begin{equation}
    D_{\text{KL}}(P_{\text{player}} \| P_{\text{engine}}) = \sum_{i}
    P_{\text{player}}(m_i) \log_2 \frac{P_{\text{player}}(m_i)}
    {P_{\text{engine}}(m_i)}.
\end{equation}

\begin{remark}[Information Overlap with the Delta Model]
\label{rem:overlap}
Since $P_{\text{player}}$ is one-hot, the KL divergence simplifies to
$D_{\text{KL}} = -\log_2 Q(m_{\text{chosen}})$, which is the negative
log-likelihood of the chosen move under the engine's policy. This is
informationally related to $\Delta$: both measure engine alignment in
different spaces (probability space vs.\ winrate space). These metrics
are \emph{not independent} signals, and practitioners should account for
this correlation to avoid overconfident conclusions.
\end{remark}

\subsection{Cross-Entropy}

\begin{equation}
    H(P, Q) = -\sum_{i} P(m_i) \log_2 Q(m_i).
\end{equation}
Since $D_{\text{KL}}(P \| Q) = H(P, Q) - H(P)$, cross-entropy provides
an unnormalized measure of engine alignment.


% =============================================================================
% 7. TEMPORAL ANALYSIS
% =============================================================================
\section{Temporal Analysis and Switch-Point Detection}
\label{sec:temporal}

\subsection{Moving Average Smoothing}

\begin{equation}
    \bar{\Delta}_t = \frac{1}{\min(t, W)} \sum_{j=\max(1,\, t-W+1)}^{t}
    \Delta_j.
\end{equation}

\subsection{Switch-Point Detection}

A switch point is detected at index $t$ if:
\begin{equation}
    |\mu_{\text{before}} - \mu_{\text{after}}| \geq \theta_{\text{switch}},
\end{equation}
where $\mu_{\text{before}}$ and $\mu_{\text{after}}$ are local means of the
smoothed sequence. A positive change (deltas decreased) indicates suspicious
improvement.

\subsection{Trend and Volatility}

Trend is estimated via linear regression slope $\hat{\beta}$ (Equation~%
\eqref{eq:trend_slope}); volatility via the coefficient of variation
$V = \sigma_\Delta / \bar{\Delta}$.

\begin{equation}
    \hat{\beta} = \frac{\sum_{i=1}^{n}(i - \bar{i})(\Delta_i - \bar{\Delta})}
    {\sum_{i=1}^{n}(i - \bar{i})^2}.
    \label{eq:trend_slope}
\end{equation}


% =============================================================================
% 8. MULTI-GAME PROFILE ANALYSIS
% =============================================================================
\section{Multi-Game Profile Analysis}
\label{sec:profile}

A single game provides insufficient statistical power for reliable detection,
as demonstrated in Section~\ref{sec:experiments}. This section develops a
formal framework for aggregating evidence across $m$ games in a player's
profile history.

\subsection{Profile Dataset Structure}

Let $\mathcal{G} = \{g_1, g_2, \ldots, g_m\}$ denote a player's profile
history of $m$ games, where game $g_j$ yields estimated rate
$\hat{\lambda}_j$ and observed mean delta $\bar{\Delta}_j$. The profile
dataset treats each game as a sample unit, yielding the profile statistics:
\begin{align}
    \bar{\lambda} &= \frac{1}{m} \sum_{j=1}^{m} \hat{\lambda}_j,
    \label{eq:profile_mean}\\
    s_\lambda &= \sqrt{\frac{1}{m-1} \sum_{j=1}^{m}
    (\hat{\lambda}_j - \bar{\lambda})^2}.
    \label{eq:profile_std}
\end{align}

The profile mean $\bar{\lambda}$ serves as an empirical estimate of the
player's true skill level, while $s_\lambda$ quantifies inter-game variability.

\subsection{Confidence Interval for the Baseline Parameter}
\label{sec:profile_ci}

Under the Central Limit Theorem (applicable for $m \geq 30$, or $t$-distribution
for $m < 30$), the $100(1-\alpha)\%$ confidence interval for the player's
true rate $\lambda_{\text{true}}$ is:
\begin{equation}
    \bar{\lambda} - t_{\alpha/2,\, m-1} \cdot \frac{s_\lambda}{\sqrt{m}}
    \;\leq\; \lambda_{\text{true}} \;\leq\;
    \bar{\lambda} + t_{\alpha/2,\, m-1} \cdot \frac{s_\lambda}{\sqrt{m}},
    \label{eq:profile_ci}
\end{equation}
where $t_{\alpha/2,\, m-1}$ is the upper $\alpha/2$ critical value of the
$t$-distribution with $m-1$ degrees of freedom.

\subsection{Anomaly Detection via Hypothesis Test}

When analyzing a suspect game with rate estimate $\hat{\lambda}_{\text{suspect}}$,
we formulate a one-sided test:
\begin{align}
    H_0 &: \lambda_{\text{suspect}} = \bar{\lambda}, \\
    H_1 &: \lambda_{\text{suspect}} > \bar{\lambda}.
\end{align}

The test statistic is:
\begin{equation}
    Z = \frac{\hat{\lambda}_{\text{suspect}} - \bar{\lambda}}
    {s_\lambda / \sqrt{m}}.
    \label{eq:profile_test}
\end{equation}

We reject $H_0$ at level $\alpha$ if $Z > z_\alpha$ (or $Z > t_{\alpha,\, m-1}$
for small $m$). This provides a formally calibrated decision rule whose
false-positive rate is controlled at $\alpha$, as opposed to the ad-hoc
posterior thresholds in Equation~\eqref{eq:classification}.

\subsection{Sample Size Requirements}
\label{sec:sample_size}

We derive formal sample size requirements for the profile analysis from
confidence-interval theory~\citep{montgomery2010applied}.

\subsubsection{Number of Games Required}

To estimate $\lambda_{\text{true}}$ with absolute margin $E$ at
$100(1-\alpha)\%$ confidence:
\begin{equation}
    m \geq \left(\frac{z_{\alpha/2} \cdot s_\lambda}{E}\right)^2.
    \label{eq:n_games}
\end{equation}

Using the pilot estimate $s_\lambda \approx 12$ obtained from preliminary
experiments (Section~\ref{sec:experiments}):

\begin{table}[H]
\centering
\caption{Required number of games $m$ for estimating $\lambda_{\text{true}}$
at 95\% confidence ($z_{0.025} = 1.96$, $s_\lambda = 12$).}
\label{tab:n_games}
\begin{tabular}{ccc}
\toprule
Margin $E$ & Formula & Required $m$ \\
\midrule
$\pm 10$ & $\lceil(1.96 \times 12 / 10)^2\rceil$ & $6$ \\
$\pm 5$  & $\lceil(1.96 \times 12 / 5)^2\rceil$  & $22$ \\
$\pm 3$  & $\lceil(1.96 \times 12 / 3)^2\rceil$  & $62$ \\
\bottomrule
\end{tabular}
\end{table}

\subsubsection{Number of Moves per Game}

For estimating the near-optimal proportion $p$ (a binomial parameter) with
margin $E$ at 95\% confidence, using the pilot estimate $\hat{p} \approx 0.87$:
\begin{equation}
    n_{\text{moves}} \geq \left(\frac{z_{\alpha/2}}{E}\right)^2
    \hat{p}(1 - \hat{p}).
    \label{eq:n_moves}
\end{equation}

\begin{table}[H]
\centering
\caption{Required moves per game for estimating near-optimal proportion
at 95\% confidence ($\hat{p} = 0.87$).}
\label{tab:n_moves}
\begin{tabular}{ccc}
\toprule
Margin $E$ & Formula & Required $n$ \\
\midrule
$\pm 0.10$ & $\lceil(19.6)^2 \times 0.113\rceil$ & $44$ \\
$\pm 0.05$ & $\lceil(39.2)^2 \times 0.113\rceil$ & $174$ \\
$\pm 0.025$ & $\lceil(78.4)^2 \times 0.113\rceil$ & $694$ \\
\bottomrule
\end{tabular}
\end{table}

For $\lambda$-based inference, the MLE variance of the Exponential rate
parameter satisfies $\text{SE}(\hat{\lambda}) = \lambda / \sqrt{n}$, yielding
the relative-error requirement:
\begin{equation}
    n \geq \left(\frac{z_{\alpha/2}}{E_{\text{rel}}}\right)^2,
    \label{eq:n_lambda}
\end{equation}
where $E_{\text{rel}}$ is the desired relative error. For $E_{\text{rel}} = 0.20$
and 95\% confidence: $n \geq (1.96/0.20)^2 = 97$ moves.

\subsubsection{Practical Recommendations}

\begin{table}[H]
\centering
\caption{Practical sample size guidelines.}
\label{tab:practical}
\begin{tabular}{lll}
\toprule
Objective & Minimum & Recommended \\
\midrule
Confidence score $\geq 80\%$ (single game) & $30$ moves & $44$ moves \\
Proportion estimate $\pm 10\%$ & $44$ moves & $100$ moves \\
$\lambda$ estimate $\pm 20\%$ & $97$ moves & $200$ moves \\
EM calibration (mixture model) & $200$ moves & $500$ moves \\
Profile baseline ($\pm 10$ in $\lambda$) & $6$ games & $20$ games \\
Profile baseline ($\pm 5$ in $\lambda$) & $22$ games & $30$ games \\
\bottomrule
\end{tabular}
\end{table}

\subsection{Stationarity Assumption and Rolling Window}

The profile analysis rests on the assumption that the player's true skill level
is approximately \emph{stationary} over the observation window.

\begin{remark}[Stationarity Assumption]
\label{rem:stationarity}
The profile mean $\bar{\lambda}$ is a valid estimator of $\lambda_{\text{true}}$
only if the player's skill level does not change systematically over the
analysis window. This assumption may be violated in two scenarios:
(i)~genuine skill improvement, which causes $\hat{\lambda}_j$ to increase
over time due to legitimate learning; and (ii)~onset of engine use at a
specific historical game, which causes a structural break in the series
$\hat{\lambda}_1, \ldots, \hat{\lambda}_m$.
\end{remark}

To address this, we recommend a \emph{rolling window} strategy: only the
$m_{\text{window}}$ most recent games are used to estimate the baseline,
with $m_{\text{window}} \in [20, 50]$ as a practical range. The temporal
switch-point detection of Section~\ref{sec:temporal} can be applied to the
game-level series $\hat{\lambda}_1, \ldots, \hat{\lambda}_m$ to detect
structural breaks in player skill, which may themselves constitute
suspicious patterns.


% =============================================================================
% 9. SYSTEM ARCHITECTURE
% =============================================================================
\section{System Architecture}
\label{sec:architecture}

\subsection{Overview}

The detection system consists of three layers:
\begin{enumerate}
    \item \textbf{Engine Layer}: An external Gomoku engine (e.g., Rapfi)
          communicates via the Gomocup protocol. For each position, the engine
          returns the top $k$ candidate moves with winrate evaluations.
    \item \textbf{Analysis Layer}: The models of Sections~\ref{sec:framework}--
          \ref{sec:temporal} compute per-move and per-game metrics.
    \item \textbf{Classification Layer}: The Bayesian posterior, LRT, temperature
          score, and profile analysis are aggregated into the final decision.
\end{enumerate}

\subsection{Analysis Pipeline}

For each move: (1) send position to engine; (2) receive top $k$ evaluations;
(3) compute $\Delta$, $C$, and $A$; (4) perform tempered Bayesian update;
(5) compute information-theoretic metrics; (6) accumulate temporal statistics.

Profile analysis runs asynchronously, aggregating per-game statistics across
the player's history and updating the baseline CI (Equation~\eqref{eq:profile_ci}).


% =============================================================================
% 10. CLASSIFICATION
% =============================================================================
\section{Classification Decision Framework}
\label{sec:classification}

\subsection{Multi-Signal Aggregation}

Classification integrates three primary signals:
\begin{enumerate}
    \item Bayesian posterior $P(H_1 \mid D)$ under stratified $H_0$.
    \item Log-likelihood ratio $\log \Lambda$ from the formal LRT.
    \item Boltzmann temperature score $S_{\tau} = 1/\hat{\tau}$.
\end{enumerate}

The primary rule is:
\begin{equation}
    \text{Classification} = \begin{cases}
        \textbf{Cheater} & P(H_1 \mid D) \geq \theta_{\text{cheat}},\\
        \textbf{Suspicious} & \theta_{\text{susp}} \leq P(H_1 \mid D)
            < \theta_{\text{cheat}},\\
        \textbf{Human} & P(H_1 \mid D) < \theta_{\text{susp}},
    \end{cases}
    \label{eq:classification}
\end{equation}
with defaults $\theta_{\text{cheat}} = 0.70$ and $\theta_{\text{susp}} = 0.30$.
Agreement between the LRT rejection and the Bayesian conclusion strengthens
the classification; disagreement warrants human review.

\subsection{Profile-Enhanced Classification}

When a profile baseline is available, the single-game posterior is supplemented
by the profile test statistic $Z$ (Equation~\eqref{eq:profile_test}). The
combined evidence is evaluated as follows:
\begin{itemize}
    \item If the single-game posterior is borderline (Suspicious) but
          $Z > z_{0.05} = 1.645$, the classification is escalated to
          Suspicious--Elevated.
    \item If the single-game posterior is below $\theta_{\text{susp}}$ but
          $Z > z_{0.01} = 2.326$, the game is flagged for manual review.
    \item A confirmed Human classification requires both sub-threshold posterior
          and $Z < z_{0.05}$.
\end{itemize}

Single-game analysis without a profile baseline should be treated as
\emph{preliminary screening} only, with formal conclusions deferred until
sufficient game history is available (Table~\ref{tab:practical}).


% =============================================================================
% 11. EXPERIMENTAL OBSERVATIONS
% =============================================================================
\section{Preliminary Experimental Observations}
\label{sec:experiments}

We report preliminary experiments conducted to validate design choices and
identify failure modes. These experiments are exploratory in nature and do not
constitute a full empirical evaluation; systematic benchmarking with labeled
datasets is planned as future work.

\subsection{Experimental Setup}

A prototype implementation was developed integrating the Rapfi engine
(via the Gomocup protocol) as the evaluation oracle. Analysis was conducted
on individual games using the V2 anti-cheat analysis pipeline. Model comparison
(Exponential, Gamma, Weibull) and EM mixture fitting were enabled.

\subsection{Case Study: Professional Player (False Positive Scenario)}

\paragraph{Setup.} We analyzed two successive runs on the same game played by
a professional human player (confirmed ground truth: no engine assistance,
22 moves analyzed, $\lambda_{\text{human}} = 10$ in Run 1, $\lambda_{\text{human}}
= 50$ in Run 2).

\paragraph{Results.}

\begin{table}[H]
\centering
\caption{Summary statistics for the professional player case study.}
\label{tab:experiment1}
\begin{tabular}{lccc}
\toprule
Metric & Run 1 ($\lambda_{\text{human}} = 10$) & Run 2 ($\lambda_{\text{human}} = 50$) \\
\midrule
Analyzed moves         & 22       & 22       \\
Mean $\Delta$          & 0.8\%    & 0.8\%    \\
$\hat{\lambda}_\text{MLE}$ & 113.32 & 113.32 \\
Near-optimal \%        & 90.9\%   & 90.9\%   \\
$P(\text{Cheat})$      & 50.3\%   & 45.2\%   \\
Classification         & Suspicious & \textbf{Human} \\
Confidence             & 73.3\%   & 73.3\%   \\
$\log \Lambda$         & 22.46    & 22.46    \\
$\hat{\tau}$           & 0.0454   & 0.0454   \\
$S_\tau$               & 22.02    & 22.02    \\
\bottomrule
\end{tabular}
\end{table}

\paragraph{Observations.}

\begin{enumerate}
    \item \textbf{False positive at $\lambda_{\text{human}} = 10$:}
          Run~1 misclassified the professional player as Suspicious.
          The posterior shifted correctly to Human in Run~2 after calibrating
          $\lambda_{\text{human}}$ to 50, confirming that the null hypothesis
          parameter must reflect the player's actual skill tier.

    \item \textbf{Inconsistency between $\log \Lambda$ and posterior:}
          The log-likelihood ratio $\log \Lambda = 22.46$ is extremely large
          (corresponding to $\Lambda = e^{22.46} \approx 5.7 \times 10^9$),
          yet the posterior $P(\text{Cheat}) \approx 45\%$--$50\%$ remains
          near the decision boundary. Sensitivity analysis shows that for
          all prior values tested, the posterior exceeds 99\%. This suggests
          an implementation inconsistency between the log-LR computation and
          the Bayesian update, likely caused by the tempered likelihood
          suppressing evidence via positional complexity. This warrants
          investigation and is consistent with the observation that complex
          positions (many with $C \ll 1$) may over-discount genuine evidence.

    \item \textbf{Weibull consistently wins model selection:}
          Across all runs, the Weibull distribution achieves the lowest AIC
          and BIC with shape parameter $k \approx 0.36$, indicating heavy-tailed
          behavior (rare but large errors). However, all three models are
          rejected by the KS test ($p < 0.05$), which is expected at $n = 22$
          due to low test power. Reliable model validation requires $n \geq 50$.

    \item \textbf{EM-fitted $\lambda_g$ is engine-level:}
          The EM algorithm consistently fits a ``good play'' component with
          $\lambda_g \approx 6{,}800$--$7{,}600$ (corresponding to
          $\E[\Delta_g] \approx 0.013\%$--$0.015\%$), far exceeding the
          cheater default $\lambda_{\text{cheater}} = 50$. This reflects the
          professional player's near-perfect execution in the majority of
          moves and underscores the overlap between the high end of the human
          distribution and engine-assisted play.

    \item \textbf{Sample size insufficiency:}
          The confidence score of 73.3\% indicates borderline data sufficiency
          ($n = 22 < n_0 = 30$). Extrapolating from Section~\ref{sec:sample_size},
          reliable proportion-based inference requires $n \geq 44$ moves per game.
          Single-game analysis at $n \leq 30$ should be treated as preliminary.
\end{enumerate}

\paragraph{Implication for Framework Design.}
This case study directly motivates three design revisions incorporated
into the present version of the framework: (i) stratified null hypothesis
(Section~\ref{sec:stratification}), (ii) profile-based baseline estimation
(Section~\ref{sec:profile}), and (iii) increased data-sufficiency threshold
$n_0$ in the confidence formula (Section~\ref{sec:bayesian}).

\subsection{Case Study: Opponent-Difficulty Confound}
\label{sec:exp_opp}

\paragraph{Setup.} We analyzed a complete game (21 moves of Black analyzed,
confirmed ground truth: no engine assistance) to investigate the effect of
opponent weakness on the complexity score. The opponent was a significantly
weaker player who lost decisively by move~33.

\paragraph{Results.}

\begin{table}[H]
\centering
\caption{Summary statistics for the opponent-difficulty case study.}
\label{tab:experiment2}
\begin{tabular}{lc}
\toprule
Metric & Value \\
\midrule
Analyzed moves              & 21      \\
Mean $\Delta$               & 3.1\%   \\
$\hat{\lambda}_\text{MLE}$  & 31.31   \\
Near-optimal \%             & 66.7\%  \\
$P(\text{Cheat})$           & 0.0\%   \\
Classification              & \textbf{Human} \\
Confidence                  & 70.0\%  \\
$\log \Lambda$              & $-4.48$ \\
$\hat{\tau}$                & 0.1262  \\
$S_\tau$                    & 7.92    \\
\bottomrule
\end{tabular}
\end{table}

\paragraph{Observations.}

\begin{enumerate}
    \item \textbf{Correct classification despite opponent weakness:}
          The framework correctly classified this game as Human
          ($P(\text{Cheat}) = 0.0\%$), driven primarily by two
          large-delta moves: Move~\#23 ($\Delta = 19.7\%$) and
          Move~\#15 ($\Delta = 17.5\%$). These blunders are characteristic
          of human play and strongly favor $H_0$.

    \item \textbf{Opponent-difficulty confound identified:}
          Inspection of the move-level log reveals a qualitative pattern
          directly motivating Section~\ref{sec:opp_complexity}. From
          Move~\#13 onward, the opponent's best available winrate dropped
          below $0.3\%$ (Quality: Winning), yet multiple positions in the
          range Move~\#13--\#33 were assigned raw complexity $C \approx
          0.47$--$0.49$ based on candidate variance alone.
          Without $C_{\text{adj}}$, the nine near-optimal moves in the
          winning phase would contribute substantial spurious evidence
          toward $H_1$ in games with lower $\Delta$ overall.

    \item \textbf{Trivial Position Override activates late:}
          The Override (Proposition~\ref{prop:trivial}) activated only at
          Move~\#35 onward, once $w^*$ exceeded $0.999$. Between Move~\#13
          and Move~\#35, a gap of eleven analyzed moves, the opponent had
          effectively lost but the Override was not triggered. This confirms
          the need for the Opponent Difficulty Factor as a complementary
          mechanism that operates on \emph{opponent} winrate rather than
          \emph{player} winrate.

    \item \textbf{Log LR is negative:}
          $\log \Lambda = -4.48$ indicates that the data are $e^{4.48}
          \approx 88$ times \emph{more likely} under $H_0$ than $H_1$,
          consistent with correct Human classification. This contrasts
          sharply with the professional player case study ($\log \Lambda =
          22.46$), demonstrating that the presence of genuine blunders
          ($\Delta \geq 15\%$) is a powerful discriminating signal.

    \item \textbf{EM identifies human-like blunder component:}
          The EM algorithm fitted $\lambda_g = 9334$ (near-perfect moves)
          and $\lambda_b = 10.69$ ($\E[\Delta_b] = 9.4\%$, genuine blunders),
          with mixing weight $\pi = 0.666$. The lower $\pi$ relative to the
          professional player case ($\pi = 0.862$) quantifies the higher
          blunder rate and is consistent with the ground-truth classification.
\end{enumerate}

\paragraph{Implication for Framework Design.}
This case study confirms that the Opponent Difficulty Factor
(Section~\ref{sec:opp_complexity}) is necessary as a complement to the
Trivial Position Override. The Override handles the extreme case of
$w^* > 0.95$ from the player's perspective; the Opponent Difficulty Factor
handles the intermediate case where the opponent is functionally defeated
but the player's winrate has not yet reached the Override threshold.
Together they provide comprehensive suppression of opponent-induced
false-positive inflation.


% =============================================================================
% 12. DISCUSSION
% =============================================================================
\section{Discussion}
\label{sec:discussion}

\subsection{Strengths}

\begin{enumerate}
    \item \textbf{Mathematical rigor.} Every component has a clear probabilistic
          interpretation grounded in established statistical theory.
    \item \textbf{Complexity awareness.} The framework reduces false positives
          from forced sequences and false negatives from selective cheating.
    \item \textbf{Skill stratification.} The stratified null hypothesis
          addresses the fundamental skill-confounding problem, enabling
          application across the full spectrum of player abilities.
    \item \textbf{Profile analysis.} Multi-game aggregation provides
          statistically sound, sample-size-justified conclusions that
          single-game analysis cannot support.
    \item \textbf{Opponent-difficulty correction.} The Opponent Difficulty
          Factor prevents false-positive inflation when the analyzed player
          faces a weak opponent, addressing a source of bias absent from
          prior frameworks.
    \item \textbf{Modularity and transparency.} Each component can be
          independently validated, updated, or replaced.
\end{enumerate}

\subsection{Limitations}

\begin{enumerate}
    \item \textbf{Engine dependency.} Analysis quality is bounded by the
          engine's evaluation accuracy. Positions where the engine errs
          produce misleading delta values.
    \item \textbf{Parameter sensitivity.} Default parameters are heuristically
          chosen. EM calibration and sensitivity analysis partially address
          this, but full calibration requires large labeled datasets.
    \item \textbf{Independence assumption.} Conditional independence of deltas
          may be violated when strategic plans span multiple moves. An HMM
          extension (Section~\ref{sec:conclusion}) relaxes this.
    \item \textbf{Selective cheating.} Sparse engine use (e.g., only at critical
          moves) may evade detection thresholds; the temporal model partially
          addresses this.
    \item \textbf{Information overlap.} KL divergence and winrate delta are not
          fully independent (Remark~\ref{rem:overlap}); over-reliance on
          correlated metrics may produce overconfident conclusions.
    \item \textbf{Stationarity.} Profile analysis assumes approximately
          stationary skill levels (Remark~\ref{rem:stationarity}). Structural
          breaks from genuine improvement may mimic the onset of cheating.
    \item \textbf{Opponent difficulty calibration.} The reference threshold
          $w_{\text{ref}} = 0.40$ and the minimum floor $\alpha_{\min} = 0.10$
          in the Opponent Difficulty Factor (Equation~\eqref{eq:c_adj}) are
          heuristically chosen. Empirical calibration from labeled game data
          is required; in particular, $w_{\text{ref}}$ may need to be
          game-phase dependent (e.g., lower in the opening than in the endgame).
    \item \textbf{Tier boundary ambiguity.} The overlap between
          $T_4$ (professional) and $H_1$ (engine-assisted) distributions
          represents a fundamental statistical limit: it may not be possible
          to reliably distinguish top professional play from engine assistance
          in a small number of games.
\end{enumerate}

\subsection{Ethical Considerations}

Anti-cheating systems carry inherent risks of false accusations. We advocate:
\begin{enumerate}
    \item Statistical evidence should inform human reviewers, not replace
          human judgment.
    \item Formal action should require high posterior probability, high
          confidence (sufficient data), and profile-level corroboration.
    \item Players should have access to the evidence and methodology used.
    \item The system should be audited regularly for skill-level and
          demographic biases.
\end{enumerate}


% =============================================================================
% 13. CONCLUSION
% =============================================================================
\section{Conclusion and Future Work}
\label{sec:conclusion}

We have presented a comprehensive Bayesian statistical framework for detecting
engine-assisted play in online Gomoku competitions. The framework integrates
parametric move-quality modeling, context-aware complexity weighting, tempered
Bayesian updates, information-theoretic divergence measures, temporal
switch-point detection, player skill stratification, and multi-game profile
analysis into a coherent probabilistic system.

Preliminary experiments reveal a fundamental \emph{skill-confounding problem}:
without stratification, professional players are incorrectly flagged at rates
comparable to engine-assisted players. The profile-based framework developed
in Section~\ref{sec:profile} provides a statistically principled resolution,
with sample size requirements derived from confidence-interval theory.
Additionally, we identify and formalize the \emph{opponent-difficulty confound}:
positions made trivially easy by a weak opponent generate spurious evidence
toward $H_1$ under one-sided analysis. The Opponent Difficulty Factor
(Section~\ref{sec:opp_complexity}) corrects for this bias by weighting
evidential contribution by the opponent's competitive viability.

\subsection{Future Work}

\begin{itemize}
    \item \textbf{Hidden Markov Model for selective cheating:}
          A two-state HMM with states $\{\text{Honest}, \text{Cheating}\}$
          and emission distributions $F_{z_t}$:
          \begin{equation}
              z_t \mid z_{t-1} \sim \text{Categorical}(A_{z_{t-1}}),
              \quad \Delta_t \mid z_t \sim F_{z_t},
          \end{equation}
          where $A$ is the transition matrix. This subsumes switch-point
          detection as a special case~\citep{rabiner1989tutorial}.

    \item \textbf{Empirical calibration:} Fitting all model parameters via
          EM, MCMC, or variational inference on large labeled datasets. KS
          and Anderson--Darling tests should validate distributional
          assumptions. This includes calibration of $w_{\text{ref}}$ and
          $\alpha_{\min}$ in the Opponent Difficulty Factor across different
          game phases and competitive contexts.

    \item \textbf{Systematic empirical evaluation:} Benchmarking on three
          scenario types---confirmed human games, engine self-play, and
          suspected games---with full ROC curve, AUC, precision--recall,
          and confusion matrix reporting.

    \item \textbf{Tier boundary calibration:} Determining empirically where
          the $T_4$/$H_1$ boundary lies using paired professional-play and
          engine-play datasets.

    \item \textbf{Dirichlet Process deployment:} Replacing the fixed two-component
          mixture with the DPMM (Section~\ref{sec:dpm}) in production.

    \item \textbf{Adversarial robustness:} Analyzing resilience against adaptive
          cheaters who intentionally introduce errors.

    \item \textbf{Generalization:} Adapting the framework to Go, Othello,
          Connect6, and other abstract strategy games.

    \item \textbf{Timing integration:} Incorporating move timing as an additional
          signal via a multi-modal Bayesian network.
\end{itemize}


% =============================================================================
% REFERENCES
% =============================================================================
\bibliographystyle{plainnat}

\begin{thebibliography}{99}

\bibitem[Regan and Haworth, 2011]{regan2011intrinsic}
Regan, K.~W. and Haworth, G.~M. (2011).
\newblock Intrinsic chess ratings.
\newblock In \emph{Proceedings of the AAAI Conference on Artificial
Intelligence}, volume~25, pages 834--839.

\bibitem[Guid and Bratko, 2006]{guid2006computer}
Guid, M. and Bratko, I. (2006).
\newblock Computer analysis of world chess champions.
\newblock \emph{ICGA Journal}, 29(2):65--73.

\bibitem[Berger and Sellke, 1987]{berger1987testing}
Berger, J.~O. and Sellke, T. (1987).
\newblock Testing a point null hypothesis: The irreconcilability of {P} values
and evidence.
\newblock \emph{Journal of the American Statistical Association},
82(397):112--122.

\bibitem[Gr{\"u}nwald and van Ommen, 2017]{grunwald2017inconsistency}
Gr{\"u}nwald, P. and van Ommen, T. (2017).
\newblock Inconsistency of {B}ayesian inference for misspecified linear models,
and a proposal for repairing it.
\newblock \emph{Bayesian Analysis}, 12(4):1069--1103.

\bibitem[Cover and Thomas, 2006]{cover2006elements}
Cover, T.~M. and Thomas, J.~A. (2006).
\newblock \emph{Elements of Information Theory}.
\newblock Wiley-Interscience, 2nd edition.

\bibitem[Bishop, 2006]{bishop2006pattern}
Bishop, C.~M. (2006).
\newblock \emph{Pattern Recognition and Machine Learning}.
\newblock Springer.

\bibitem[Dempster et~al., 1977]{dempster1977maximum}
Dempster, A.~P., Laird, N.~M., and Rubin, D.~B. (1977).
\newblock Maximum likelihood from incomplete data via the {EM} algorithm.
\newblock \emph{Journal of the Royal Statistical Society: Series B},
39(1):1--38.

\bibitem[Ferguson, 1973]{ferguson1973bayesian}
Ferguson, T.~S. (1973).
\newblock A {B}ayesian analysis of some nonparametric problems.
\newblock \emph{The Annals of Statistics}, 1(2):209--230.

\bibitem[Neal, 2000]{neal2000markov}
Neal, R.~M. (2000).
\newblock Markov chain sampling methods for {D}irichlet process mixture models.
\newblock \emph{Journal of Computational and Graphical Statistics},
9(2):249--265.

\bibitem[Blei and Jordan, 2006]{blei2006variational}
Blei, D.~M. and Jordan, M.~I. (2006).
\newblock Variational inference for {D}irichlet process mixtures.
\newblock \emph{Bayesian Analysis}, 1(1):121--143.

\bibitem[Luce, 1959]{luce1959individual}
Luce, R.~D. (1959).
\newblock \emph{Individual Choice Behavior: A Theoretical Analysis}.
\newblock John Wiley \& Sons.

\bibitem[Neyman and Pearson, 1933]{neyman1933most}
Neyman, J. and Pearson, E.~S. (1933).
\newblock On the problem of the most efficient tests of statistical hypotheses.
\newblock \emph{Philosophical Transactions of the Royal Society of London,
Series A}, 231:289--337.

\bibitem[Rabiner, 1989]{rabiner1989tutorial}
Rabiner, L.~R. (1989).
\newblock A tutorial on hidden {M}arkov models and selected applications in
speech recognition.
\newblock \emph{Proceedings of the IEEE}, 77(2):257--286.

\bibitem[Montgomery and Runger, 2010]{montgomery2010applied}
Montgomery, D.~C. and Runger, G.~C. (2010).
\newblock \emph{Applied Statistics and Probability for Engineers}.
\newblock John Wiley \& Sons, 5th edition.

\end{thebibliography}

\end{document}
